% ===== PRÉAMBULE CANONIQUE — à coller VERBATIM en tête de CHAQUE .tex =====
\documentclass[11pt,a4paper]{article}
\usepackage[utf8]{inputenc}\usepackage[T1]{fontenc}\usepackage{lmodern}
\usepackage[french]{babel}
\usepackage{amsmath,amssymb,mathtools}\usepackage{icomma}
\usepackage[version=4]{mhchem}          % \ce{...} : équations & formules
\usepackage{siunitx}\sisetup{locale=FR} % \qty{}{}, \si{}, \ang{}
\usepackage{chemfig}                    % \chemfig{...} : structures développées
\usepackage{xcolor}\usepackage{colortbl}
\usepackage{tikz}\usepackage{pgfplots}\pgfplotsset{compat=1.18}
\usepackage[most]{tcolorbox}\usepackage{enumitem}\usepackage{array,booktabs}
\usepackage[a4paper,margin=2.2cm]{geometry}
\usetikzlibrary{arrows.meta,calc,angles,quotes,positioning,patterns,backgrounds,shapes.geometric}
\definecolor{primary}{RGB}{222,49,99}\definecolor{primarydark}{RGB}{150,22,55}
\definecolor{defcol}{RGB}{0,114,178}\definecolor{methcol}{RGB}{52,92,148}
\definecolor{excol}{RGB}{0,135,90}\definecolor{attncol}{RGB}{200,40,55}
\tcbset{boxbase/.style={enhanced,breakable,boxrule=0.9pt,arc=2.5pt,
  left=3mm,right=3mm,top=2.3mm,bottom=2.3mm,fonttitle=\bfseries,coltitle=white}}
% --- boîtes TUTORIEL (noms EXACTS, ne pas renommer) ---
\newtcolorbox{cle}[1][]{boxbase,colback=defcol!6,colframe=defcol,
  title={\textsc{À retenir}\ifx\relax#1\relax\relax\else\textmd{ : \itshape #1}\fi}}
\newtcolorbox{methode}[1][]{boxbase,colback=methcol!7,colframe=methcol,sharp corners,
  title={\textsc{Méthode}\ifx\relax#1\relax\relax\else\textmd{ --- \itshape #1}\fi}}
\newtcolorbox{exemplebox}[1][]{boxbase,colback=excol!5,colframe=excol,
  title={\textsc{Exemple}\ifx\relax#1\relax\relax\else\textmd{ : \itshape #1}\fi}}
\newtcolorbox{piege}[1][]{boxbase,colback=attncol!6,colframe=attncol,
  title={\textsc{Piège}\ifx\relax#1\relax\relax\else\textmd{ : \itshape #1}\fi}}
\newtcolorbox{remarque}[1][]{boxbase,colback=black!4,colframe=black!45,
  title={\textsc{Remarque}\ifx\relax#1\relax\relax\else\textmd{ : \itshape #1}\fi}}
% --- boîtes EXERCICE / CORRIGÉ (noms EXACTS, auto-comptées) ---
\newtcolorbox[auto counter]{exo}[1][]{boxbase,colback=primary!4,colframe=primary,
  title={\textsc{Exercice}~\thetcbcounter\ifx\relax#1\relax\relax\else\textmd{ --- \itshape #1}\fi}}
\newtcolorbox[auto counter]{corrige}[1][]{boxbase,colback=excol!4,colframe=excol,
  title={\textsc{Corrigé}~\thetcbcounter\ifx\relax#1\relax\relax\else\textmd{ --- \itshape #1}\fi}}
\newcommand{\R}{\mathbb{R}}\newcommand{\N}{\mathbb{N}}\newcommand{\Z}{\mathbb{Z}}\newcommand{\ds}{\displaystyle}
% (un \usepackage{fancyhdr}+en-tête est possible ; il est ignoré côté web)
% =========================================================================
\begin{document}

\begin{center}{\large\bfseries\color{primarydark} Thème 3 — Corrigé Facile}\end{center}

\begin{corrige}[concentration molaire]
\begin{enumerate}[label=\alph*)]
  \item $C = \dfrac{0{,}5}{2} = \SI{0.25}{\mole\per\litre}$.
  \item $C = \dfrac{0{,}2}{0{,}5} = \SI{0.40}{\mole\per\litre}$.
  \item $C = \dfrac{3}{4} = \SI{0.75}{\mole\per\litre}$.
\end{enumerate}
\end{corrige}

\begin{corrige}[concentration massique]
\begin{enumerate}[label=\alph*)]
  \item $C_m = \dfrac{20}{2} = \SI{10.0}{\gram\per\litre}$.
  \item $C_m = \dfrac{5}{0{,}5} = \SI{10.0}{\gram\per\litre}$.
  \item $C_m = \dfrac{45}{5} = \SI{9.00}{\gram\per\litre}$.
\end{enumerate}
\end{corrige}

\begin{corrige}[quantité de matière dissoute]
\begin{enumerate}[label=\alph*)]
  \item $n = 0{,}2\times 3 = \SI{0.6}{\mole}$.
  \item $n = 1{,}5\times 0{,}2 = \SI{0.3}{\mole}$.
  \item $\SI{500}{\milli\litre} = \SI{0.5}{\litre}$ : $n = 0{,}1\times 0{,}5 = \SI{0.05}{\mole}$.
\end{enumerate}
\end{corrige}

\begin{corrige}[passer d'une concentration à l'autre]
\begin{enumerate}[label=\alph*)]
  \item $C_m = C\,M = 0{,}1\times 40 = \SI{4.00}{\gram\per\litre}$.
  \item $C_m = 0{,}5\times 58{,}5 = \SI{29.25}{\gram\per\litre}$.
  \item $C = \dfrac{C_m}{M} = \dfrac{20}{40} = \SI{0.500}{\mole\per\litre}$.
\end{enumerate}
\end{corrige}

\begin{corrige}[attention aux unités de volume]
\begin{enumerate}[label=\alph*)]
  \item $\SI{250}{\milli\litre} = \SI{0.25}{\litre}$ : $C = \dfrac{0{,}05}{0{,}25} = \SI{0.200}{\mole\per\litre}$.
  \item $\SI{500}{\milli\litre} = \SI{0.5}{\litre}$ : $C_m = \dfrac{8}{0{,}5} = \SI{16.0}{\gram\per\litre}$.
\end{enumerate}
\end{corrige}

\end{document}
